% Source MD: E:\wanru\PRML_Course\NoteMD\Week1\第3节课 机器学习通用指南：模型+损失函数+优化.md
\documentclass[main.tex]{subfiles}

\graphicspath{{\subfix{figures/}}}

\begin{document}

\chapter{机器学习通用指南：模型、损失函数与优化}
\label{chap:w01l03-model-loss-optimization}

\begin{examplebox}
这节课读起来更像一份排错手册。实验结果不理想时，先别急着乱调模型，先分清问题卡在训练端还是测试端，再去判断是模型偏差、优化问题、过拟合，还是数据分布不一致。
\end{examplebox}

\section{机器学习任务的基本框架}
\label{sec:w01l03-basic-framework}

后面做作业或 Kaggle 型任务时，手上的资料通常会分成两块：
\begin{itemize}
  \item \textbf{Training data}：包含输入特征 $X$ 与真实标签 $Y$。
  \item \textbf{Testing data}：只有输入特征 $X$，需要模型输出预测结果。
\end{itemize}

任务长相可以差很多，但训练主线经常很像。无论是语音辨识、影像辨识还是机器翻译，大多都逃不开这三步：
\begin{enumerate}
  \item 写出带未知参数的函数 $f_\theta(x)$；
  \item 定义损失函数 $L(\theta)$；
  \item 解决优化问题，找出最佳参数 $\theta^\star$。
\end{enumerate}

\begin{fact}{三步训练流程}{w01l03-three-steps}
大多数监督学习任务都可以抽象为模型设计、损失定义与优化求解三步。
\end{fact}

统一的目标写成
\begin{equation}
  \theta^\star = \argmin_{\theta} L\paren{\theta}.
  \label{eq:w01l03-optimization}
\end{equation}

这条式子正好把前两课接起来：在线性模型里，$\theta$ 可能只有 $\paren{w,b}$；换成神经网络，$\theta$ 就会一路扩成矩阵 $W$、偏置向量 $b$、输出层权重 $c$ 等一大串参数。
外观虽然越来越复杂，但训练想做的事没变，都是把参数往能让 loss 更小的方向推。

\section{训练误差过大：Model Bias 与 Optimization Issue}
\label{sec:w01l03-training-loss}

Kaggle 分数不好看时，最容易犯的错就是直接盯着 leaderboard 一通乱改。
课程给的顺序更稳：先回头看训练集上的 loss，再决定问题到底出在哪一层。

\subsection{Model Bias}
\label{sec:w01l03-model-bias}

如果模型太简单，可表示的函数集合本来就很小，那么就算已经把这类模型里最好的那一个找出来，loss 还是可能很高。
这种时候，不是你没认真训练，而是模型一开始就不具备足够的表达能力。

常见解法包括：
\begin{itemize}
  \item 增加输入特征；
  \item 增加神经元数量或层数；
  \item 直接换成更有弹性的模型家族。
\end{itemize}

\subsection{Optimization Issue}
\label{sec:w01l03-optimization-issue}

另一类情况就不是模型太弱，而是参数没被顺利找到。
课程里给了一个很实用的判断法：如果一个更深、更大的模型，训练集 loss 反而比浅层模型还高，那通常不像模型偏差，更像是优化过程出了问题。

\begin{warningbox}
深层模型的表达能力比浅层模型更强。如果训练 loss 反而更差，问题通常出在优化过程，而不是模型不够大。
\end{warningbox}

\section{测试误差过大：Overfitting}
\label{sec:w01l03-overfitting}

如果训练 loss 已经压得很低，测试 loss 却还是不好看，那排查方向就该转向过拟合。

\begin{definition}{过拟合（Overfitting）}{w01l03-overfitting}
模型在训练资料上拟合得很好，却带不去没见过的资料，这种情形就叫过拟合（overfitting）。
\end{definition}

课程里给的对策，基本都围着“别让模型把训练资料记死”这件事展开：
\begin{enumerate}
  \item \textbf{增加训练资料}：这是最有效的方法。
  \item \textbf{Data augmentation}：在不破坏任务语义的前提下，对资料进行合理扩增。
  \item \textbf{限制模型复杂度}：例如减少层数、减少参数、减少特征。
  \item \textbf{使用额外正则手段}：例如 early stopping、regularization、dropout。
\end{enumerate}

\begin{examplebox}
对影像资料做左右翻转通常合理，但上下颠倒往往会破坏真实语义。数据扩增不是随便造数据，而是要遵守任务本身的结构。
\end{examplebox}

\subsection{Bias-Complexity Trade-off}
\label{sec:w01l03-bias-complexity}

这一段可以把它想成两条曲线在分家。
模型复杂度刚开始提高时，训练 loss 往下走，测试 loss 通常也会跟着改善；但复杂到某个程度以后，训练 loss 还能继续降，测试 loss 却可能开始回头上升。
从函数集合的角度看，这不难理解：模型越复杂，可表示的函数越多，训练集上的最优 loss 往往只会更低；可测试集面对的是没见过的资料，一旦模型开始连训练资料里的偶然噪声都记住，泛化就会变差。
这就是课程里说的 bias-complexity trade-off：模型太简单不行，复杂到失控也不行。

\section{模型选择与交叉验证}
\label{sec:w01l03-model-selection}

课程特别提醒，不要直接拿 Kaggle public leaderboard 当模型选择器。
如果你每次都根据 public 分数来改模型，时间一久，其实就在对 public set 本身过拟合。

更稳妥的做法是从训练资料里切出 validation set：
\begin{itemize}
  \item 用 training set 训练模型；
  \item 用 validation set 比较多个候选模型；
  \item 选出最佳模型后，再用全部训练资料重训一次。
\end{itemize}

如果担心一次切分不稳定，还可以用 $N$-fold cross validation：
\begin{enumerate}
  \item 把资料切成 $N$ 份；
  \item 每次拿其中 1 份做 validation，其余做 training；
  \item 轮流跑完 $N$ 次；
  \item 比较各模型在 $N$ 次验证上的平均表现。
\end{enumerate}

\section{Mismatch：训练分布与测试分布不同}
\label{sec:w01l03-mismatch}

\begin{definition}{分布不一致（Mismatch）}{w01l03-mismatch}
如果 training data 和 testing data 根本不是从同一个数据分布来的，那么就算模型在训练集上表现不错，测试时也可能整段垮掉。这种现象称为分布不一致（mismatch）。
\end{definition}

分布不一致和一般过拟合不太一样。这里的问题不是模型太会记，而是训练资料和测试资料本来就处在不同情境里。
像拿疫情期间的流量模式去预测疫情后的年份，这类数据分布偏移就会很明显。

\end{document}
